\section{Essentially self-dual inducing blocks}
\label{sec:essential-selfdual-blocks}

\begin{theorem}[D-BI-09: the degree-twenty-four gate]
\label{D-BI-09}
Suppose a finite-cover Weil sixfold source of exact imaginary quadratic
character field $K$ has a period map of rank nine and nonsplit actual
polarization.  Suppose its complex coefficient representation is
$\operatorname{Ind}_H^G\tau$, where $\tau$ is irreducible and primitive
of degree $k\ge9$.  If $\tau$ is essentially self-dual, then
\[
                         k\ge24.                         \tag{BI.17}
\]
Here essentially self-dual means that there is a nondegenerate complex
bilinear form $B$ and a character $\alpha$ of the finite image $J$ of
$\tau$ such that
\[
              g^{\mathsf T}Bg=\alpha(g)B\qquad(g\in J).
                                                               \tag{BI.18}
\]
The form may be symmetric or alternating.  The inducing character field
$L=\mathbf Q(\chi_\tau)$ need not equal $K$ and no model of $\tau$
over $K$ is assumed.  In particular, the conclusion applies to every
primitive inducing block that is real up to scalar matrices.

For a factorwise tensor decomposition $\tau=A\otimes B_0$ with
$2\le a=\dim A\le b=\dim B_0$ and $k=ab$, the stronger necessary
inequality is
\[
 (a-1)b\ge 9a-6+\epsilon_k,
 \qquad \epsilon_k=\begin{cases}0&k\text{ even},\\1&k\text{ odd}.
 \end{cases}                                                 \tag{BI.19}
\]
The tensor factors in this assertion may be projective representations.
\end{theorem}
\begin{proof}
Pass to the actual intermediate connected $J$-cover $C_0\to X_0$
as in D-BI-01.  Its $\tau$-multiplicity variation is the original
complex multiplicity variation, by Frobenius reciprocity and quotient
pullback.  In particular its cohomological multiplicity has dimension
six, and its period differential has the same rank as the original
$K$-factor: the conjugate filtration is its polarization annihilator.
If $h=g(X_0)$ and $D$ is the sum of the moving dimensions of the
genuine local inertias, then
\[
                         D=6-2k(h-1).
\]
For $h=1$, the marked-base moduli has dimension at most six; for
$h\ge2$ this formula is negative because $k\ge9$.  A finite cover
of a fixed marked base has only finitely many topological possibilities.
Consequently rank nine forces $h=0$ and at least twelve genuine
branch points.  Write their nonidentity inertia matrices as
$g_1,\ldots,g_r$.  They generate $J$, have product one, and satisfy
\[
 r\ge12,\qquad \sum_{j=1}^r\operatorname{rank}(1-g_j)=2k+6.
                                                               \tag{BI.20}
\]
This is a factorization through the actual intermediate Hurwitz family;
it is not a subtraction of dimensions of unrelated parameter spaces.

Induction of characters gives $K\subset L$.  If $\alpha$ were
trivial, (BI.18) would make $\tau$ self-dual.  The character of a
finite self-dual complex representation is real, contradicting the
inclusion of the imaginary quadratic field $K$ in $L$.  Thus $\alpha$
is nontrivial.  Since the $g_j$ generate and their product is one,
at least two of them have $\alpha(g_j)\ne1$.

The transpose of $B$ has the same equivariance, so Schur's lemma
gives $B^{\mathsf T}=B$ or $-B$ after rescaling; in fact the sign
already follows by transposing twice.  If $\alpha(g)\ne1$ and
$x,y\in\ker(g-1)$, (BI.18) yields
$B(x,y)=\alpha(g)B(x,y)$, hence $B(x,y)=0$.  The fixed subspace
is totally isotropic for a nondegenerate symmetric or alternating
form and therefore has dimension at most $\lfloor k/2\rfloor$.
The two distinguished inertias consequently have moving rank at least
\[
                            p=\lceil k/2\rceil.             \tag{BI.21}
\]

First suppose $J$ is tensor indecomposable in the convention of
\cite[paragraph before Theorem~7.1]{BI:GS}.  If it contains an element
of moving rank two, the characteristic-zero classification
\cite[Theorem~7.1 and Remark~(3), pp.~548--549]{BI:GS} gives either
the scalar deleted alternating/symmetric case or the degree-ten
$U_5(2)$ case.  The first is excluded by D-BI-08, which explicitly
allows the larger inducing field and the full subgroup range
$\mathfrak A_{k+1}\subset J\subset\mathbf C^*\mathfrak S_{k+1}$.
In the second case the finite image normalizes a nondegenerate
alternating form, as in D-BI-02 and
\cite[Table~II]{BI:GT}.  Its multiplier character must again be
nontrivial, since otherwise $L$ would be real.  Two inertias move at
least five dimensions.  Every remaining nonidentity inertia moves
at least two: a symplectic semisimple transformation has positive
moving rank at least two, and a nontrivial similitude moves at least
five.  Thus the total is at least $2\cdot5+10\cdot2=30$, exceeding
$2k+6=26$.  This also excludes the exceptional case.

There are no reflections either.  Indeed, after averaging a positive
Hermitian form, a reflection's moving line is perpendicular to its
fixed hyperplane.  Irreducibility gives two conjugate reflections
$u,v$ with distinct moving lines.  If $uvx=x$, then
$(v-1)x=(u^{-1}-1)x$ lies in the intersection of those two lines,
which is zero.  Hence
$\ker(uv-1)=\ker(u-1)\cap\ker(v-1)$, of codimension two.
Their product would be the already excluded bireflection.

Every other genuine inertia now moves at least three dimensions.
Combining this with (BI.20)--(BI.21) gives
\[
              2k+6\ge 2\lceil k/2\rceil+30
                         =k+\epsilon_k+30.
                                                               \tag{BI.22}
\]
Therefore $k\ge24+\epsilon_k$, proving the assertion in this case.
Groups that permute tensor factors but do not preserve a factorwise
decomposition fall within the cited tensor-indecomposable convention;
they are not omitted here.

It remains to handle a factorwise tensor decomposition of dimensions
$2\le a\le b$.  For any tensor matrix, the fixed space has dimension
at most the largest eigenspace multiplicity in the first factor
times $b$.  Thus a matrix whose first projective factor is nontrivial
moves at least $b$ dimensions.  If only the second projective factor
is nontrivial, it moves at least $a$; a nonidentity scalar matrix
moves all $k$ dimensions.  Every nonidentity inertia therefore moves
at least $a$.

At least three inertia matrices have nontrivial first projective
factor.  Otherwise their product-one relation makes the group
generated by the first projective factors cyclic.  A cyclic projective
group is represented by powers of one diagonalizable matrix up to
scalars and has an invariant line.  Tensoring that line with the
second factor contradicts irreducibility.  This argument does not
require the projective factors to admit linear lifts as representations.

Choose two of the inertias in (BI.21).  At least one of those three
nontrivial first-projective-factor inertias remains outside this pair.
The two selected inertias contribute at least $2p$, that third one
at least $b$, and at least nine others contribute at least $a$ each.
Consequently
\[
 2ab+6\ge 2\lceil ab/2\rceil+b+9a
                 =ab+\epsilon_k+b+9a,
\]
which is precisely (BI.19).  If $ab<24$, then $a\le4$.
For $a=2$, (BI.19) requires $b\ge12$; for $a=3$ it requires
$b\ge11$; for $a=4$ it requires $b\ge10$.  Each contradicts
$a\le b$ and $ab<24$.  This completes the proof.
\end{proof}

\begin{proposition}[D-BI-10: the multiplier field and odd degrees]
\label{D-BI-10}
Let $V$ be an absolutely irreducible primitive representation defined
over an imaginary quadratic field $K$, with exact character field
$K$, and suppose it is essentially self-dual.  Its similitude
character is unique and takes values in the roots of unity of $K$.
If $K\ne\mathbf Q(i),\mathbf Q(\sqrt{-3})$, then $\dim V$ is even.
In particular, in these fields an essentially self-dual primitive
coefficient representation satisfying D-BI-09 has even degree at
least twenty-four.  This last field assertion concerns $V$ itself;
it is not imposed on a larger-field inducing representation.
\end{proposition}
\begin{proof}
First a primitive irreducible finite representation has no nontrivial
linear self-twist.  If an invertible intertwiner $T$ satisfies
$T\rho(g)=\psi(g)\rho(g)T$, where $\psi$ has order $d>1$,
then $T^d$ is scalar.  Rescale $T$ so that $T^d=1$.
Its eigenspaces are permuted by the full image $\mu_d$ of $\psi$.
All $d$ eigenvalues occur, and their nonzero eigenspaces form a
transitively permuted direct-sum decomposition, contrary to primitivity.

Two nondegenerate equivariant bilinear forms with multipliers
$\alpha_1,\alpha_2$ give, by their quotient, a self-twist with
character $\alpha_2/\alpha_1$.  The preceding argument therefore
shows that the multipliers are identical and the forms proportional.
Choose an algebraic form $B$; this is possible because its defining
linear equations have algebraic coefficients and its multiplier has
finite image.  Apply any automorphism of $\overline{\mathbf Q}$
fixing $K$ to (BI.18).  The representation matrices are unchanged,
so uniqueness implies $\alpha^\gamma=\alpha$.  Thus $\alpha$
is $K$-valued.  The character is nontrivial because otherwise the
exact character field would be real.

For the specified fields the only roots of unity are $\pm1$.
The determinant character also takes values in $\{\pm1\}$.
Taking determinants in (BI.18) yields
\[
                   \det(\rho(g))^2=\alpha(g)^{\dim V}.
\]
The left side is one.  A nontrivial $\alpha$ takes the value $-1$,
which makes the right side minus one if $\dim V$ is odd.  This
contradiction proves the parity assertion.
\end{proof}
